\phantomsection
\section*{Глава 3. Разработка веб-приложения}
\addcontentsline{toc}{section}{Глава 3. Разработка веб-приложения}
\stepcounter{section}
\label{sec:chapter3}

\providecommand{\tabref}[1]{\hyperref[#1]{таблице~\ref*{#1}}}
\providecommand{\figref}[1]{\hyperref[#1]{рисунке~\ref*{#1}}}
\providecommand{\appref}[1]{\hyperref[#1]{приложении~\ref*{#1}}}

Целью третьего раздела является описание практической реализации веб-приложения для учета, хранения и синхронизации конфигурационных и диагностических данных ПЛК LogicBox. В первой главе была определена предметная область, а во второй главе спроектированы роли пользователей, функции приложения, модель данных, архитектура интерфейсов и механизм взаимодействия с локальным клиентом. В данной главе рассматривается, как указанные проектные решения были реализованы в программном коде.

При разработке особое внимание уделялось не только набору функций, но и однозначности поведения системы. Для промышленной предметной области важно, чтобы пользователь не мог случайно нарушить иерархию объектов эксплуатации, удалить связанные данные, загрузить файл в неправильный контекст или получить доступ к данным вне своей области ответственности. Поэтому большинство ограничений реализовано не только на уровне интерфейса, но и на уровне серверной логики и базы данных (БД). Такой подход позволяет снизить зависимость безопасности от клиентской части приложения и делает систему более устойчивой к ручной подмене запросов.


\subsection{Разработка базовой структуры приложения и модуля аутентификации}

Серверная часть приложения реализована на языке Go с использованием веб-фреймворка Fiber (фреймворк для обработки HTTP-запросов и построения веб-приложений)~\cite{golang_official,fiber_official}. В качестве системы управления базами данных используется PostgreSQL (реляционная система управления базами данных), так как она поддерживает внешние ключи, ограничения целостности, транзакции, триггеры и расширенные проверки на уровне схемы данных~\cite{postgresql_official}. Клиентская часть реализована с использованием HTML (HyperText Markup Language --- язык гипертекстовой разметки), CSS (Cascading Style Sheets --- каскадные таблицы стилей) и JavaScript (язык программирования для клиентской логики браузера)~\cite{mdn_html,mdn_css,mdn_js}.

Структура проекта разделена на несколько логических частей. Точка входа находится в файле \texttt{main.go} каталога \texttt{cmd/web}; в ней выполняется загрузка конфигурации, подключение к БД, применение миграций, подготовка шаблонизатора и регистрация маршрутов. Серверная логика расположена в каталоге приложения, настройки вынесены в отдельный конфигурационный модуль, работа с БД отделена от файлового хранилища. Пользовательские шаблоны размещены в каталоге \texttt{web/templates}, статические файлы интерфейса --- в \texttt{web/static}, а SQL-миграции (SQL, Structured Query Language --- язык структурированных запросов) --- в каталоге \texttt{migrations}.

Основные элементы структуры приложения представлены в \tabref{tab:chapter3_structure}.

\begin{center}
    \small
    \captionof{table}{Основные элементы программной структуры приложения}
    \label{tab:chapter3_structure}
    \begin{tabularx}{\textwidth}{|>{\raggedright\arraybackslash}p{0.34\textwidth}|>{\raggedright\arraybackslash}X|}
        \hline
        \textbf{Элемент проекта} & \textbf{Назначение} \\ \hline
        \texttt{cmd/web/}\newline\texttt{main.go} & Запуск приложения, подключение к БД, применение миграций и регистрация HTTP-маршрутов. \\ \hline
        \texttt{internal/app/}\newline\texttt{server.go} & Основная серверная логика: маршруты, проверка прав, работа с объектами, ПЛК, конфигурациями, файлами, архивом и администрированием. \\ \hline
        \texttt{internal/database/}\newline\texttt{database.go} & Подключение к PostgreSQL, последовательное применение SQL-миграций и первичное заполнение демонстрационных данных. \\ \hline
        \texttt{internal/storage/}\newline\texttt{local.go} & Сохранение загруженных файлов, проверка размера и расширения, формирование безопасного имени и вычисление контрольной суммы. \\ \hline
        \texttt{web/templates} & HTML-шаблоны страниц: панель, объекты, контроллеры, конфигурации, диагностика, файлы, архив, синхронизация и администрирование. \\ \hline
        \texttt{migrations} & Последовательные изменения схемы БД: таблицы, ограничения, индексы и триггеры для контроля целостности. \\ \hline
    \end{tabularx}
    \normalsize
\end{center}

Запуск приложения выполняется из файла \texttt{main.go}. В нем сначала загружаются переменные окружения, затем устанавливается соединение с PostgreSQL, после чего при включенном режиме автоматического применения миграций (\texttt{AUTO\_MIGRATE}) обновляется схема БД. Такой подход удобен для развертывания, поскольку структура БД автоматически приводится к актуальному состоянию при запуске приложения.

\begin{samepage}
\begin{lstlisting}[caption={Фрагмент инициализации веб-приложения}, label={lst:chapter3_main_start}]
cfg := config.Load()
db, err := database.Open(ctx, cfg.DatabaseURL)
if err != nil {
    log.Fatalf("connect database: %v", err)
}

if cfg.AutoMigrate {
    if err := database.ApplySchema(ctx, db, "migrations"); err != nil {
        log.Fatalf("migrate: %v", err)
    }
}

engine := html.New("./web/templates", ".html")
fiberApp := fiber.New(fiber.Config{
    Views: engine,
    BodyLimit: cfg.MaxUploadSizeMB * 1024 * 1024,
})
\end{lstlisting}
\end{samepage}

Как видно из листинга~\ref{lst:chapter3_main_start}, ограничение размера тела запроса устанавливается на уровне конфигурации приложения. Это важно для загрузки файлов, поскольку система работает с конфигурациями, журналами, архивами и отчетами, но не должна принимать неконтролируемо большие файлы.

Для аутентификации пользователей реализована форма входа. Пользователь вводит логин и пароль, после чего сервер ищет учетную запись в таблице \texttt{users}, проверяет ее статус и сравнивает введенный пароль с bcrypt-хэшем. Bcrypt представляет собой алгоритм хэширования паролей, предназначенный для безопасного хранения учетных данных. В БД не хранится исходный пароль, что снижает риск компрометации учетных записей при несанкционированном доступе к данным.

После успешного входа в серверной сессии сохраняется идентификатор пользователя. Сессия хранится в cookie-файле браузера с параметром \texttt{HttpOnly}, что ограничивает доступ к ней из клиентского JavaScript. Дополнительно для всех защищенных маршрутов используется проверка активного статуса учетной записи: если пользователь заблокирован администратором, его существующая сессия уничтожается при следующем обращении к системе.

\begin{samepage}
\begin{lstlisting}[caption={Проверка активной сессии пользователя}, label={lst:chapter3_auth_check}]
func (s *Server) requireAuth(c *fiber.Ctx) error {
    u, err := s.currentUser(c)
    if err != nil || u == nil {
        return c.Redirect("/login")
    }
    if u.Status != "active" {
        sess, _ := s.Store.Get(c)
        _ = sess.Destroy()
        return c.Redirect("/login")
    }
    c.Locals("user", u)
    return c.Next()
}
\end{lstlisting}
\end{samepage}

Таким образом, аутентификация в приложении реализована как обязательный вход в систему с последующей проверкой состояния учетной записи при каждом обращении к закрытым страницам.


\subsection{Реализация ролевой модели и системы проверки доступа}

В приложении используется комбинированная модель авторизации. Авторизация означает проверку прав пользователя на выполнение конкретного действия после успешной аутентификации. Проверка выполняется по двум направлениям: по роли пользователя и по закрепленной области объектов эксплуатации.

В системе реализованы следующие роли:
\begin{itemize}
    \item администратор --- управление пользователями, областями доступа, локальными клиентами и всеми данными системы;
    \item инженер --- работа с закрепленными объектами, контроллерами, файлами, диагностикой и конфигурациями;
    \item ведущий инженер --- выполнение инженерных операций и подтверждение актуальности конфигураций;
    \item наблюдатель --- просмотр разрешенных данных без права изменения.
\end{itemize}

При этом одной роли недостаточно для доступа ко всем данным. Для обычных пользователей дополнительно применяется таблица доступа пользователей к объектам \texttt{user\_object\_access}; в ней задаются объект эксплуатации и уровень доступа. Уровни доступа имеют следующие значения: \texttt{read} --- просмотр, \texttt{write} --- изменение, \texttt{verify} --- изменение и подтверждение конфигураций. Если доступ назначен к верхнему объекту иерархии, например к организации или участку, то пользователь получает доступ также к дочерним объектам. Для этого в SQL-запросах применяется рекурсивная выборка.

\begin{samepage}
\begin{lstlisting}[caption={Проверка доступа к объекту с учетом уровня доступа}, label={lst:chapter3_access_check}]
func (s *Server) canAccessObject(ctx context.Context, u *User,
    objectID int64, write bool) bool {
    if u == nil {
        return false
    }
    if u.RoleName == "admin" {
        return true
    }
    if write && u.RoleName == "observer" {
        return false
    }
    // Далее проверяется доступ к объекту или его родителю
    // через рекурсивную выборку object_ancestors.
}
\end{lstlisting}
\end{samepage}

Данная проверка применяется не только при отображении кнопок в интерфейсе, но и в обработчиках серверных запросов. Это принципиально важно: если поле или кнопка скрыты в браузере, пользователь технически все равно может отправить запрос вручную. Поэтому сервер самостоятельно проверяет, имеет ли пользователь право на изменение конкретного объекта, контроллера, конфигурации или файла.

На уровне маршрутов реализовано дополнительное разделение. Административные страницы вынесены в группу \texttt{/admin} и доступны только пользователям с ролью \texttt{admin}. Операции просмотра и инженерного редактирования дополнительно проверяются внутри соответствующих обработчиков.

\begin{samepage}
\begin{lstlisting}[caption={Регистрация основных маршрутов приложения}, label={lst:chapter3_routes}]
private := app.Group("", s.requireAuth)
private.Get("/objects", s.objectsList)
private.Get("/controllers", s.controllersList)
private.Get("/configurations", s.configurationsList)
private.Get("/archive", s.archiveList)
private.Get("/sync/sessions", s.syncSessionsList)

admin := private.Group("/admin", s.requireRole("admin"))
admin.Get("/users", s.adminUsers)
admin.Get("/access", s.adminAccess)
admin.Get("/sync-clients", s.adminSyncClients)
admin.Get("/audit", s.adminAudit)
\end{lstlisting}
\end{samepage}

Кроме проверки доступа, в приложении реализована защита от CSRF (Cross-Site Request Forgery --- подделка межсайтового запроса). Для всех обычных POST-запросов сервер проверяет защитный токен формы. Исключение сделано только для серверного API (Application Programming Interface --- программный интерфейс приложения) синхронизации, так как он использует отдельную аутентификацию по Bearer-токену.

\begin{samepage}
\begin{lstlisting}[caption={Проверка CSRF-токена для изменяющих запросов}, label={lst:chapter3_csrf}]
func (s *Server) csrfGuard(c *fiber.Ctx) error {
    method := strings.ToUpper(c.Method())
    if method != fiber.MethodPost && method != fiber.MethodPut &&
        method != fiber.MethodPatch && method != fiber.MethodDelete {
        return c.Next()
    }
    if strings.HasPrefix(c.Path(), "/api/v1/") {
        return c.Next()
    }
    // Далее выполняется сравнение токена из сессии
    // и токена, переданного формой.
}
\end{lstlisting}
\end{samepage}

Таким образом, в системе используется не декоративная защита на уровне интерфейса, а полноценная серверная проверка прав и контекста выполнения операции.


\subsection{Реализация иерархической модели объектов эксплуатации}

Одним из центральных решений приложения является строгая иерархия объектов эксплуатации. В системе выделены четыре типа объекта: организация, участок, стенд и шкаф. Иерархия имеет фиксированный порядок:

\begin{center}
    \textit{организация $\rightarrow$ участок $\rightarrow$ стенд $\rightarrow$ шкаф}.
\end{center}

Такое ограничение делает модель данных однозначной. Организация является верхним уровнем и не имеет родителя. Участок может находиться только внутри организации, стенд --- только внутри участка, а шкаф --- только внутри стенда. ПЛК может быть привязан только к шкафу. В текущей версии системы не используется промежуточный статус \texttt{limited}, так как ограничение доступа решается не статусом объекта, а таблицей назначений доступа. Для объектов оставлены два понятных состояния: \texttt{active} --- рабочий объект, \texttt{detached} --- объект, отвязанный от родителя перед удалением.

На общей странице создания объекта пользователь может выбрать тип и родительский объект самостоятельно. Однако при создании объекта из карточки родителя форма заранее фиксирует допустимый тип. Например, из карточки организации создается только участок, из карточки участка --- только стенд, из карточки стенда --- только шкаф. Из карточки шкафа дочерние объекты не создаются, так как шкаф является нижним уровнем иерархии. Это снижает вероятность ошибки пользователя и делает процесс создания элементов более понятным.

При этом фиксация типа в интерфейсе не является единственной защитой. На сервере тип и родитель дополнительно проверяются по контексту запроса. Если пользователь попытается изменить скрытое поле формы вручную, сервер отклонит действие как нарушение иерархии.

\begin{samepage}
\begin{lstlisting}[caption={Проверка размещения объекта в иерархии}, label={lst:chapter3_object_placement}]
func (s *Server) validateObjectPlacement(ctx context.Context,
    u *User, objectID int64, typ string, parentID any) (string, any, error) {

    typ = normalizeObjectType(typ)
    if typ == "" {
        return "", nil, fmt.Errorf("выберите корректный тип объекта")
    }
    // Проверяются родитель, тип родителя, права пользователя
    // и отсутствие циклической связи.
    return typ, parent, nil
}
\end{lstlisting}
\end{samepage}

Дополнительные ограничения реализованы в БД через триггеры PostgreSQL. Триггер является процедурой, которая автоматически выполняется при изменении данных в таблице. В данном случае триггеры запрещают некорректную иерархию даже в том случае, если изменение выполняется не через интерфейс приложения, а прямым SQL-запросом.

\begin{samepage}
\begin{lstlisting}[caption={Ограничение допустимых статусов объекта на уровне БД}, label={lst:chapter3_object_status_sql}]
ALTER TABLE exploitation_objects
ADD CONSTRAINT chk_exploitation_objects_status
CHECK (status IN ('active','detached'));
\end{lstlisting}
\end{samepage}

Для безопасного удаления объектов реализована двухэтапная логика. Объект можно удалить только в том случае, если он пустой и не связан с родителем. Если объект является дочерним, его сначала нужно отвязать от родителя, после чего он получает статус \texttt{detached}. Только после этого возможно удаление. При расчете пустоты объекта учитываются дочерние объекты, активные ПЛК, локальные клиенты и сессии синхронизации. Архивные ПЛК не блокируют отвязку шкафа, поскольку они уже выведены из рабочей иерархии и доступны через раздел архива.

Такой подход исключает ситуацию, при которой пользователь удаляет часть иерархии и оставляет связанные данные в неопределенном состоянии. В интерфейсе также отображается блок безопасности удаления, где показано количество дочерних объектов, ПЛК, локальных клиентов и сессий синхронизации. Это делает причину блокировки удаления понятной для пользователя.


\subsection{Реализация CRUD-интерфейсов для объектов, ПЛК, конфигураций и диагностики}

В приложении реализованы CRUD-операции (Create, Read, Update, Delete --- создание, чтение, изменение и удаление) для основных сущностей системы. К ним относятся объекты эксплуатации, ПЛК, конфигурации, диагностические записи, файлы-вложения, пользователи, назначения доступа, локальные клиенты и сессии синхронизации.

Главная особенность реализации заключается в том, что рабочие списки и архив разделены. Обычные страницы контроллеров, конфигураций, вложений и диагностики отображают только активные данные. Заархивированные ПЛК не занимают шкаф, не выводятся в рабочих списках и не мешают расчету пустоты объекта. При этом связанные с архивным ПЛК конфигурации, диагностические записи и файлы не теряются: они доступны через карточку архивного ПЛК.

Состав основных интерфейсов приложения представлен в \tabref{tab:chapter3_interfaces}.

\begin{center}
    \small
    \captionof{table}{Основные пользовательские интерфейсы приложения}
    \label{tab:chapter3_interfaces}
    \begin{tabularx}{\textwidth}{|>{\raggedright\arraybackslash}p{0.30\textwidth}|>{\raggedright\arraybackslash}X|}
        \hline
        \textbf{Раздел} & \textbf{Назначение} \\ \hline
        Панель & Отображение сводных показателей по объектам, контроллерам, конфигурациям, диагностике и последним сессиям синхронизации. \\ \hline
        Объекты & Просмотр иерархии объектов эксплуатации, создание дочерних объектов, редактирование, отвязка и удаление пустых объектов. \\ \hline
        Контроллеры & Учет ПЛК LogicBox, привязка к шкафу, редактирование сведений, загрузка конфигураций, архивация и безопасное удаление. \\ \hline
        Конфигурации & Просмотр версий конфигураций, назначение текущей версии, подтверждение ведущим инженером или администратором, скачивание файлов. \\ \hline
        Диагностика & Отображение диагностических записей с фильтрацией по доступным пользователю объектам и ПЛК. \\ \hline
        Файлы & Загрузка и скачивание вложений, связанных с ПЛК, конфигурациями, диагностикой или сессиями синхронизации. \\ \hline
        Архив & Просмотр архивных ПЛК, связанных конфигураций, диагностических записей и файлов без включения этих данных в рабочие списки. \\ \hline
        Администрирование & Управление пользователями, областями доступа, локальными клиентами и журналом аудита. \\ \hline
    \end{tabularx}
    \normalsize
\end{center}

При создании ПЛК сервер проверяет, что выбранный объект является доступным шкафом. Данная проверка выполняется независимо от того, пришел ли запрос из общей формы или из карточки шкафа. Это защищает систему от подмены \texttt{object\_id} в форме.

\begin{samepage}
\begin{lstlisting}[caption={Проверка привязки ПЛК только к шкафу}, label={lst:chapter3_controller_cabinet}]
func (s *Server) isAvailableCabinetObject(ctx context.Context,
    objectID int64) bool {
    var ok bool
    err := s.DB.QueryRowContext(ctx, `
        SELECT EXISTS(
            SELECT 1 FROM exploitation_objects
            WHERE id=$1
              AND type='cabinet'
              AND status='active'
        )`, objectID).Scan(&ok)
    return err == nil && ok
}
\end{lstlisting}
\end{samepage}

Для конфигураций реализованы операции загрузки, просмотра, назначения текущей версии, подтверждения и удаления. Назначение текущей версии выполняется в транзакции. Транзакция представляет собой группу операций с БД, которые выполняются как единое целое: если одна операция завершается ошибкой, изменения откатываются. Это важно при переключении текущей конфигурации, поскольку сначала необходимо снять признак текущей версии у остальных записей, а затем установить его у выбранной.

\begin{samepage}
\begin{lstlisting}[caption={Назначение текущей конфигурации в транзакции}, label={lst:chapter3_current_config}]
tx, err := s.DB.BeginTx(c.Context(), nil)
defer tx.Rollback()

_, err = tx.ExecContext(c.Context(),
    `UPDATE configurations SET is_current=false
      WHERE controller_id=$1 AND is_deleted=false`, controllerID)

_, err = tx.ExecContext(c.Context(),
    `UPDATE configurations SET is_current=true
      WHERE id=$1 AND is_deleted=false`, id)

err = tx.Commit()
\end{lstlisting}
\end{samepage}

Подтверждение конфигураций доступно только администратору и ведущему инженеру. Это отражает реальную инженерную практику: не каждая загруженная конфигурация автоматически считается проверенной и пригодной для дальнейшего использования. В системе разделены понятия текущей версии и подтвержденной версии. Текущая версия показывает, какая конфигурация является актуальной для ПЛК, а подтверждение фиксирует факт проверки ответственным специалистом.

Диагностические записи в текущей версии в основном поступают через сессии синхронизации. Они содержат уровень, тип события, код, сообщение, время события и связь с конкретным ПЛК. В списках диагностики отображаются только записи, доступные пользователю по его области ответственности. Это позволяет использовать общий журнал без раскрытия данных по чужим объектам.


\subsection{Реализация загрузки, скачивания и безопасного хранения файлов}

Работа с файлами является одной из ключевых функций системы, так как конфигурации, журналы, отчеты и архивы часто существуют в виде отдельных файлов. Для их хранения реализован локальный файловый слой. При загрузке файл сохраняется в каталог, заданный переменной \texttt{UPLOAD\_DIR}. Внутри этого каталога создается структура по категории, году и месяцу. Такой подход упрощает обслуживание хранилища и предотвращает накопление всех файлов в одной директории.

Перед сохранением файла выполняются проверки:
\begin{itemize}
    \item проверяется наличие файла в запросе;
    \item проверяется максимальный размер файла;
    \item проверяется расширение файла;
    \item исходное имя очищается от опасных символов;
    \item для файла вычисляется контрольная сумма SHA-256;
    \item фактическое имя на диске создается случайным идентификатором.
\end{itemize}

Контрольная сумма SHA-256 (Secure Hash Algorithm 256-bit --- криптографическая хэш-функция с длиной результата 256 бит) используется для фиксации содержимого файла и дальнейшей проверки неизменности. В БД сохраняются исходное имя, MIME-тип (Multipurpose Internet Mail Extensions --- тип содержимого файла), размер, путь хранения и контрольная сумма.

\begin{samepage}
\begin{lstlisting}[caption={Сохранение файла в управляемом каталоге}, label={lst:chapter3_storage_save}]
func (s *LocalStorage) Save(fileHeader *multipart.FileHeader,
    category string) (SavedFile, error) {

    if fileHeader.Size > s.MaxUploadBytes {
        return SavedFile{}, fmt.Errorf("file is too large")
    }
    ext := strings.ToLower(filepath.Ext(fileHeader.Filename))
    if !allowedExt(ext) {
        return SavedFile{}, fmt.Errorf("file extension is not allowed")
    }
    // Далее создается каталог и файл сохраняется
    // с расчетом контрольной суммы SHA-256.
}
\end{lstlisting}
\end{samepage}

Особое внимание уделено скачиванию и удалению файлов. Перед выдачей файла сервер проверяет, что путь к нему находится внутри управляемого каталога загрузок. Это защищает систему от ситуации, когда в БД ошибочно или злонамеренно записан путь к файлу вне каталога приложения. Если путь не проходит проверку, приложение отказывает в скачивании или удалении.

\begin{samepage}
\begin{lstlisting}[caption={Проверка принадлежности пути каталогу загрузок}, label={lst:chapter3_storage_path}]
func (s *Server) isManagedStoragePath(path string) bool {
    rootAbs, err := filepath.Abs(s.Cfg.UploadDir)
    if err != nil {
        return false
    }
    pathAbs, err := filepath.Abs(path)
    if err != nil {
        return false
    }
    return pathAbs == rootAbs ||
        strings.HasPrefix(pathAbs, rootAbs+string(os.PathSeparator))
}
\end{lstlisting}
\end{samepage}

При ошибке записи информации о файле в БД приложение удаляет уже сохраненный физический файл. Это предотвращает появление бесхозных файлов, не связанных ни с одной записью системы. При удалении конфигурации или архивного ПЛК также выполняется удаление связанных физических файлов после успешного изменения данных в БД.


\subsection{Реализация архивации ПЛК и безопасного жизненного цикла данных}

В процессе разработки была реализована отдельная логика жизненного цикла ПЛК. Она основана на различии между пустым ПЛК и ПЛК, у которого уже есть связанные данные. Если ПЛК не имеет конфигураций, диагностических записей и файлов, его можно удалить напрямую. Если же к ПЛК привязаны конфигурации, диагностика или файлы, сначала он переводится в архив.

Архивация не удаляет связанные данные. Она переводит ПЛК в состояние \texttt{archived}, фиксирует дату архивации и сохраняет идентификатор последнего места установки. После этого ПЛК не отображается в обычных списках и не занимает шкаф. Это важно для иерархии: шкаф, содержащий только архивные ПЛК, считается пустым и может быть отвязан или удален по общим правилам. При этом архивные конфигурации, диагностические записи и файлы остаются доступными через карточку архивного ПЛК.

Разделение рабочих и архивных данных устраняет противоречие между двумя требованиями. С одной стороны, архивные ПЛК не должны мешать текущей эксплуатации и отображаться в обычных списках. С другой стороны, исторические данные нельзя терять, так как они могут понадобиться для анализа изменений, расследования неисправностей или восстановления информации о предыдущем состоянии оборудования.

Основные правила жизненного цикла ПЛК представлены в \tabref{tab:chapter3_controller_lifecycle}.

\begin{center}
    \small
    \captionof{table}{Правила жизненного цикла ПЛК в системе}
    \label{tab:chapter3_controller_lifecycle}
    \begin{tabularx}{\textwidth}{|>{\raggedright\arraybackslash}p{0.34\textwidth}|>{\raggedright\arraybackslash}X|}
        \hline
        \textbf{Ситуация} & \textbf{Действие системы} \\ \hline
        ПЛК не имеет конфигураций, диагностических записей и файлов & ПЛК можно удалить напрямую; запись скрывается флагом удаления. \\ \hline
        ПЛК имеет связанные данные & Прямое удаление блокируется; сначала требуется архивирование. \\ \hline
        ПЛК переведен в архив & ПЛК исключается из рабочих списков и перестает занимать шкаф. \\ \hline
        Пользователь открывает карточку архивного ПЛК & Система показывает последнее место установки, конфигурации, диагностику и архивные файлы. \\ \hline
        Пользователь удаляет архивный ПЛК & Система помечает связанные данные как удаленные и удаляет связанные физические файлы из управляемого каталога. \\ \hline
    \end{tabularx}
    \normalsize
\end{center}

В карточке архивного ПЛК доступны конфигурации и файлы. Конфигурации можно открыть для просмотра, а связанные файлы можно скачать. Изменение, подтверждение и удаление архивной конфигурации заблокированы до разархивирования ПЛК. Это делает архивный режим режимом хранения и просмотра, а не активного редактирования.

На уровне БД также исправлена логика отвязки объектов: при проверке пустоты шкафа учитываются только активные ПЛК, а архивные ПЛК исключаются из расчета. Благодаря этому архив не мешает обслуживанию иерархии объектов.


\subsection{Реализация приема данных от локального клиента}

Для интеграции с локальным контуром объекта реализован серверный API приема данных. API используется локальным клиентом, который находится в сети объекта, получает сведения о ПЛК LogicBox и передает их на сервер. В текущей версии система не выполняет удаленную запись конфигураций в ПЛК и не отправляет управляющие команды. Это ограничение является осознанным, так как управляющее воздействие на промышленное оборудование требует отдельного анализа рисков.

Передача данных выполняется в формате JSON (JavaScript Object Notation --- текстовый формат обмена данными). Локальный клиент обращается к маршруту \texttt{/api/v1/sync/}\allowbreak\texttt{import} и передает Bearer-токен. Bearer-токен представляет собой секретную строку, которая передается в заголовке запроса и подтверждает право локального клиента на импорт данных. На сервере токен сравнивается с bcrypt-хэшем, сохраненным в таблице локальных клиентов синхронизации \texttt{sync\_clients}.

\begin{samepage}
\begin{lstlisting}[caption={Общая структура запроса синхронизации}, label={lst:chapter3_sync_json}]
{
  "object_id": 15,
  "client_name": "notebook-metro-vom1",
  "controllers": [
    {
      "name": "LB-VOM1-MAIN",
      "ip_address": "10.10.1.11",
      "model": "LogicBox 241",
      "status": "active"
    }
  ],
  "diagnostics": [
    {
      "controller_name": "LB-VOM1-MAIN",
      "level": "warning",
      "event_type": "voltage",
      "message": "Напряжение ниже уставки"
    }
  ]
}
\end{lstlisting}
\end{samepage}

Сервер проверяет, что локальный клиент имеет право работать именно с указанным объектом. Если в запросе передан другой идентификатор объекта (\texttt{object\_id}), импорт отклоняется. Также проверяется, что импорт контроллеров выполняется только для объекта типа \texttt{cabinet} (внутренний код типа объекта «шкаф»). После успешной проверки создается запись сессии синхронизации, обновляются или создаются сведения о контроллерах и добавляются диагностические записи. Далее в тексте используется термин «шкаф», а значение \texttt{cabinet} остается только внутренним кодом системы.

\begin{samepage}
\begin{lstlisting}[caption={Фрагмент проверки локального клиента при импорте}, label={lst:chapter3_sync_check}]
if req.ObjectID > 0 && req.ObjectID != objectID {
    return c.Status(403).JSON(fiber.Map{
        "error": "sync client is not allowed for requested object",
    })
}
if len(req.Controllers) > 0 &&
    !s.isAvailableCabinetObject(c.Context(), objectID) {
    return c.Status(400).JSON(fiber.Map{
        "error": "controllers can be imported only for cabinet",
    })
}
\end{lstlisting}
\end{samepage}

В результате каждая синхронизация становится прослеживаемой. В системе фиксируются локальный клиент, объект, время начала и завершения, IP-адрес источника, количество переданных контроллеров и диагностических записей, а также исходное краткое содержимое запроса. Это позволяет использовать журнал синхронизации для проверки поступления данных и последующего анализа.


\subsection{Реализация аналитических и административных функций}

Для удобства работы инженера и администратора в приложении реализованы аналитические страницы. Панель показывает общее количество доступных объектов, активных контроллеров, конфигураций, диагностических записей и последних сессий синхронизации. Отдельный раздел аналитики агрегирует диагностические записи по уровню, объектам и периоду времени, показывает распределение ПЛК по статусам, покрытие конфигурациями, состояние синхронизаций и потенциально устаревшие контроллеры.

Аналитика строится с учетом прав пользователя. Администратор видит данные по всей системе, а инженер или наблюдатель --- только по тем объектам, к которым ему назначен доступ. Это важно, так как даже агрегированные данные могут раскрывать сведения о чужих объектах эксплуатации.

Административный раздел включает управление пользователями, областями доступа, локальными клиентами и журналом аудита. При создании пользователя пароль сразу хэшируется. При удалении пользователя запись не удаляется физически, а переводится в статус \texttt{blocked}. Такой подход сохраняет историю аудита и предотвращает потерю связей с ранее выполненными действиями.

В журнал аудита записываются ключевые действия: вход и выход пользователя, создание и изменение объектов, загрузка конфигураций, подтверждение версий, архивация ПЛК, удаление файлов, назначение доступа и импорт сессий синхронизации. Для промышленной системы это особенно важно, поскольку позволяет установить, кто и когда выполнил действие, связанное с конфигурационными или диагностическими данными.


\subsection{Обеспечение стабильности и безопасности реализованной системы}

На этапе разработки были выделены потенциально опасные ситуации, которые могли привести к нарушению целостности данных. Для каждой такой ситуации была реализована защита на уровне серверной логики, БД или файлового хранилища. Основные меры приведены в \tabref{tab:chapter3_safety_measures}.

\begin{center}
    \small
    \captionof{table}{Основные меры обеспечения безопасности и стабильности}
    \label{tab:chapter3_safety_measures}
    \begin{tabularx}{\textwidth}{|>{\raggedright\arraybackslash}p{0.35\textwidth}|>{\raggedright\arraybackslash}X|}
        \hline
        \textbf{Риск} & \textbf{Реализованная защита} \\ \hline
        Подмена типа объекта при отправке формы & Сервер повторно проверяет допустимый тип, родителя и права пользователя; БД дополнительно проверяет иерархию триггером. \\ \hline
        Создание ПЛК не в шкафу & Сервер и БД разрешают привязку ПЛК только к объекту типа «шкаф»; во внутреннем коде этот тип обозначен как \texttt{cabinet}. \\ \hline
        Удаление объекта со связанными данными & Перед удалением проверяются дочерние объекты, активные ПЛК, локальные клиенты и сессии синхронизации. \\ \hline
        Потеря истории ПЛК при удалении & ПЛК со связанными данными сначала архивируется, а связанные конфигурации и файлы остаются доступны через архивную карточку. \\ \hline
        Доступ пользователя к чужому объекту & Все выборки и изменяющие операции проверяют роль пользователя и его область доступа. \\ \hline
        CSRF-атака на POST-формы & Для обычных изменяющих запросов проверяется CSRF-токен из серверной сессии. \\ \hline
        Скачивание файла вне каталога загрузок & Перед выдачей или удалением файла проверяется, что путь находится внутри управляемого каталога \texttt{UPLOAD\_DIR}. \\ \hline
        Оставление файла на диске после ошибки БД & Если запись в БД не создана, уже сохраненный физический файл удаляется. \\ \hline
        Использование заблокированной учетной записи & При каждом обращении к закрытым страницам проверяется статус пользователя; заблокированная сессия уничтожается. \\ \hline
    \end{tabularx}
    \normalsize
\end{center}

Особое значение имеет то, что ограничения продублированы на разных уровнях. Интерфейс помогает пользователю выполнить действие правильно, но не считается единственным механизмом защиты. Сервер проверяет все критичные параметры заново, а БД контролирует наиболее важные инварианты: допустимые типы объектов, статусы, уровни доступа, иерархию и условия безопасного удаления. Инвариант в данном контексте означает правило, которое должно оставаться истинным при любом допустимом изменении данных.

В результате система сохраняет однозначное поведение при основных сценариях работы:
\begin{itemize}
    \item объект создается только в допустимом месте иерархии;
    \item ПЛК создается только внутри шкафа;
    \item активные данные отображаются в рабочих списках;
    \item архивные данные не мешают рабочим спискам, но доступны через архив;
    \item пустой объект можно отвязать и удалить;
    \item объект со связанными активными данными удалить нельзя;
    \item пользователь видит и изменяет только данные в рамках назначенного доступа;
    \item физические файлы не выдаются и не удаляются вне управляемого каталога.
\end{itemize}

Применение таких правил делает приложение более устойчивым к ошибкам пользователя и к ручной подмене параметров запроса. Для рассматриваемой предметной области это является одним из основных требований, поскольку система работает с инженерными данными ПЛК и должна обеспечивать не только удобство учета, но и сохранность истории эксплуатации.


\subsection{Выводы по разработке веб-приложения}

В результате выполнения третьего раздела было разработано веб-приложение для централизованного учета ПЛК LogicBox, хранения конфигураций, диагностических записей, файлов и результатов синхронизации. Реализованы аутентификация пользователей, ролевая модель доступа, разграничение данных по иерархии объектов эксплуатации, CRUD-интерфейсы, файловое хранилище, архив ПЛК, журнал синхронизации, аналитические страницы и административные функции.

Ключевым результатом является не только наличие отдельных страниц и форм, но и формирование устойчивой логики жизненного цикла данных. Иерархия объектов стала фиксированной и однозначной, статусы объектов сокращены до реально используемых состояний, архивные ПЛК отделены от рабочих списков, а все опасные операции проверяются на сервере и в БД. Благодаря этому приложение можно рассматривать как единую систему сопровождения инженерных данных, а не как набор независимых таблиц и форм.

Разработанная реализация соответствует требованиям, сформулированным в первой и второй главах: она обеспечивает учет объектов и ПЛК, работу с конфигурационными и диагностическими данными, загрузку и скачивание файлов, прием данных от локального клиента, разграничение доступа пользователей и защиту целостности данных. При этом система не выполняет удаленное управление промышленным оборудованием, что соответствует принятому ограничению по снижению рисков для АСУТП.
